\item \textbf{Energía disponible en el océano: fisión vs fusión.}
Los océanos del mundo contienen un gran depósito de energía. A su grupo se le ha encomendado la tarea de determinar qué proporcionaría más energía: extrayendo el uranio en el océano para utilizarlo en reactores de fisión, o extrayendo el deuterio en el océano para utilizarlo en reactores de fusión. Divida a su grupo en dos mitades y encuentre la energía disponible en el océano desde cada fuente.
\begin{description}
    \item[Grupo (i): Combustible de fisión.] El agua de mar contiene $3.00\text{ mg}$ de uranio por metro cúbico. Alrededor del $0.700\%$ del uranio natural es el isótopo fisionable $^{235}\text{U}$.
    \item[Grupo (ii): Combustible de fusión.] De todo el hidrógeno en los océanos, $0.0300\%$ de la masa es deuterio. Dos deuterones se fusionan para formar helio en la forma $^{4}_{2}\text{He}$. Supongamos que todo el deuterio en los océanos se fusiona para formar helio.
\end{description}
Para ambos grupos, use el hecho de que la profundidad promedio del océano es de aproximadamente $4.00\text{ km}$ y el agua cubre dos tercios de la superficie de la Tierra.